% Figura: arquitectura de un VLM generativo tipo LLaVA. Adaptada de liu2023visual (Figura 1).
% Se usa en sec:marco_vlm_generativos, después del párrafo que presenta el esquema.

\begin{figure}[htbp]
    \centering
    \begin{tikzpicture}[x=1mm, y=1mm]

        % Rama visual.
        \node[figentrada, text width=18mm] (img) at (0, 0) {imagen};
        \node[figtarea, text width=22mm] (enc) at (0, 16) {codificador visual\\{\footnotesize (tipo CLIP)}};
        \node[figtarea, text width=22mm] (proj) at (0, 32) {proyector\\{\footnotesize (capa lineal)}};
        \draw[figflecha] (img) -- (enc);
        \draw[figflecha] (enc) -- (proj);

        % Rama de texto.
        \node[figentrada, text width=22mm] (instr) at (40, 0) {instrucción};
        \node[fignodo, text width=22mm, fill=figrelleno] (tok) at (40, 16) {tokenizador y\\\emph{embeddings}};
        \draw[figflecha] (instr) -- (tok);

        % Secuencia de entrada al LLM: tokens visuales como prefijo.
        \foreach \k in {0,...,3} {
            \node[draw=figazul!80!black, fill=figazul!12, minimum width=5mm, minimum height=5mm]
                (v\k) at (4+\k*6, 46) {};
        }
        \foreach \k in {0,...,3} {
            \node[draw=figtrazo, fill=figrelleno, minimum width=5mm, minimum height=5mm]
                (t\k) at (32+\k*6, 46) {};
        }
        \node[figetiqueta, anchor=south] at (13, 49.5) {tokens visuales};
        \node[figetiqueta, anchor=south] at (41, 49.5) {tokens de texto};
        \draw[figflecha] (proj.north) -- ++(0, 5) -| (13, 43.5);
        \draw[figflecha] (tok.north) -- ++(0, 21.5) -| (41, 43.5);

        % LLM y salida.
        \node[figtarea, minimum width=58mm] (llm) at (27, 64)
            {LLM de solo decodificador};
        \draw[figflecha] (27, 51) -- (llm.south);
        \node[figsalida, text width=30mm] (out) at (27, 80) {texto generado};
        \draw[figflecha] (llm) -- (out);

        \node[figetiqueta, text=fignaranja!70!black, align=left, anchor=west] at (58, 32)
            {etapa 1: se entrena\\solo el proyector};
        \node[figetiqueta, text=fignaranja!70!black, align=left, anchor=west] at (58, 64)
            {etapa 2: proyector\\y LLM};

    \end{tikzpicture}
    \caption[Arquitectura de un modelo visión--lenguaje generativo]{Arquitectura de un modelo
    visión--lenguaje generativo del tipo de LLaVA. El codificador visual extrae los rasgos de la imagen y el
    proyector los lleva a la dimensión de los \emph{embeddings} del LLM. Los tokens visuales resultantes
    entran como prefijo de la instrucción, y el LLM genera la respuesta token a token. A la derecha, qué se
    entrena en cada etapa del ajuste. Adaptada de \textcite{liu2023visual}.}
    \label{fig:marco_llava}
\end{figure}
